% ---------------------------------------------------------------------------
% Abkuerzungsverzeichnis
% Abgeglichen mit der Confluence-Seite "Abkuerzungsverzeichnis (Entwurf)"
% (Seiten-ID 101384193), Stand 05.10.2026. Aufgenommen ist nur, was in der
% Kurzfassung tatsaechlich vorkommt; ISO, JSON, SQL und UI sind deshalb entfallen.
% ---------------------------------------------------------------------------
\addchap{Abkürzungsverzeichnis}

\begin{description}[
	style=nextline,
	leftmargin=3.2cm,
	labelwidth=2.8cm,
	labelsep=0.4cm,
	font=\normalfont\bfseries,
	itemsep=0.2em]
	\item[A1 bis D7] Kennungen der 25 Testaufgaben in den Blöcken A bis D (Tabelle~\ref{tab:aufgaben})
	\item[API] Application Programming Interface, Programmierschnittstelle
	\item[API-Contract] Eigene, unveröffentlichte Schnittstellenbeschreibung des Projekts. Sie legt Datenmodell, Rechte, Zustandsautomat und Endpunkte fest, bevor sie umgesetzt werden; Verweise darauf nennen den Abschnitt, etwa §4.1.
	\item[BI] Behavioral Intention, Nutzungsabsicht (Konstrukt des TAM)
	\item[CSS] Cascading Style Sheets
	\item[FA] Funktionale Anforderung (FA-01 bis FA-44, Tabelle~\ref{tab:fa})
	\item[H1 bis H10] Die zehn Usability-Heuristiken nach Nielsen (2024) in den Bezeichnungen der Tabelle~\ref{tab:heuristiken}
	\item[HfG] Hochschule für Gestaltung Schwäbisch Gmünd
	\item[HSB] Hochschule Bremen
	\item[HTTP] Hypertext Transfer Protocol
	\item[ID] Kennung (Identifier)
	\item[NFA] Nicht-funktionale Anforderung (NFA-01 bis NFA-62, Tabelle~\ref{tab:nfa})
	\item[P1 bis P30] Kennungen der Usability-Probleme (Anhang~\ref{anh:problemliste})
	\item[PDF] Portable Document Format
	\item[PEOU] Perceived Ease of Use, wahrgenommene Benutzerfreundlichkeit (Konstrukt des TAM)
	\item[PU] Perceived Usefulness, wahrgenommene Nützlichkeit (Konstrukt des TAM)
	\item[REST] Representational State Transfer
	\item[SUS] System Usability Scale (Brooke, 1996)
	\item[TAM] Technology Acceptance Model, Technologieakzeptanzmodell (Davis, 1989)
	\item[TF] Teilfrage (TF 1 bis TF 3, Abschnitt~\ref{sec:ziele})
	\item[TN01 bis TN11] Pseudonyme der elf Teilnehmenden des Nutzertests
	\item[UC] Use Case, Anwendungsfall (UC-01 bis UC-05, Anhang~\ref{anh:user-stories})
	\item[US] User Story (US-01 bis US-42, Anhang~\ref{anh:user-stories})
	\item[W3C] World Wide Web Consortium
	\item[WCAG] Web Content Accessibility Guidelines (W3C, 2023)
\end{description}
